% Copyright (c) 2026 Ahmed Taha Bouallegui
% SPDX-License-Identifier: GPL-3.0-or-later

% =============================================================================
%  Ahmed Taha Bouallegui -- CV (version française)
%  Style : moderncv "classic" (the iconic name-banner + section-rule look)
%  Color : blue (matches personal branding)
%  Compile: pdflatex bouallegui-cv-fr.tex   (or xelatex), from this directory
%  Le contenu suit bouallegui-cv-en.tex : mêmes projets, mêmes dates, mêmes chiffres.
% =============================================================================
\documentclass[a4paper,11pt]{moderncv}
\moderncvstyle{classic}
\moderncvcolor{blue}
% Vrais glyphes pour les coordonnees et les deux adresses de profil, au lieu des
% abbreviations "in" / "gh" que le style classic utilise par defaut.
\moderncvicons{awesome}
% Le bloc des coordonnées est un peu plus serré que la taille par défaut : il partage
% la ligne d'en-tête avec le nom et la photo, et à 10pt l'icône de l'e-mail touchait
% la dernière lettre du nom.
\renewcommand*{\addressfont}{\fontsize{9}{11}\selectfont\mdseries\slshape}

% Fonts / encoding
\usepackage[utf8]{inputenc}
\usepackage{lmodern}   % scalable Type1 fonts; microtype refuses to run with bitmap ones
\usepackage{graphicx}
% Toutes les images du CV sont a cote de la source, dans img/.
\graphicspath{{img/}}
\usepackage[french]{babel}

% Puces maitrisees : babel-french redefinit itemize (tiret demi-cadratin, grande marge
% gauche) apres tweaklist de moderncv. Un environnement "list" simple passe outre.
\newenvironment{cvlist}
  {\begin{list}{\strut\textcolor{color1}{\large\rmfamily\textbullet}}{\setlength{\leftmargin}{13pt}\setlength{\labelsep}{6pt}\setlength{\itemsep}{0pt}\setlength{\topsep}{1pt}\setlength{\parsep}{0pt}\setlength{\partopsep}{0pt}}}
  {\end{list}}

% Tighten geometry so everything fits on one page
\usepackage[top=7mm,bottom=5mm,left=12mm,right=12mm]{geometry}
\setlength{\parskip}{0pt}\setlength{\itemsep}{0pt}\setlength{\partopsep}{0pt}
% Un peu plus de place pour les dates, pour qu'elles tiennent sur une ligne.
\setlength{\hintscolumnwidth}{0.195\textwidth}

% Contact details
\name{Ahmed Taha}{Bouallegui}
\phone{+33 6 68 51 84 86}
\email{ahmedtbou@tutamail.com}
\address{Villeneuve d'Ascq, France}

% La photo d'identité est la convention française ; le CV anglais n'en a pas.
% 54pt et non 60pt : le nom, les coordonnées et la photo se partagent la même ligne
% d'en-tête, et à 60pt l'icône de l'e-mail touchait la dernière lettre du nom.
\photo[54pt][0.4pt]{portrait-cv.jpg}

% Logo de l'école ou de l'entreprise, placé devant le nom de l'institution.
\newcommand{\cvlogo}[2][11pt]{\raisebox{-0.3em}{\includegraphics[height=#1]{#2}}\hspace{4pt}}

\begin{document}
\makecvtitle
% Sur papier il n'y a pas d'hypertexte, donc les trois adresses de profil sont imprimees
% en entier. Elles passent sous le filet du titre : l'adresse LinkedIn complete fait
% 246pt de large, et dans l'en-tete elle ne tiendrait a cote du nom et de la photo que
% si le nom descendait a 22pt. Le glyphe \faGlobe vient du jeu fontawesome5 deja
% charge par moderncv pour les icones sociales.
\vspace{-14pt}
\begingroup\small\color{color2}\centering
\faGlobe\hspace{4pt}bouallegui.pages.dev\hspace{1.5em}\linkedinsocialsymbol www.linkedin.com/in/ahmed-taha-bouallegui-474215306\hspace{1.5em}\githubsocialsymbol github.com/TahaBouallegui\par
\endgroup
\footnotesize

% ---------------------------------------------------------------------------
\section{Profil technique}
\cvitem[.05em]{}{Étudiant en ingénierie, spécialisé IoT et cybersécurité. Systèmes Linux, infrastructure déclarative avec Nix sur NixOS, home lab inclus qui fait tourner ses modèles en local. Stage visé en administration systèmes ou cybersécurité, 10 semaines à 4 mois, entre le 19 avril et le 31 août 2027.}

% ---------------------------------------------------------------------------
\section{Formation}
\cventry[.05em]{Depuis 2023}{M.Eng Internet of Things and Cybersecurity (prévu 2028)}{\cvlogo[13pt]{polytech-mark.png}École Polytechnique Universitaire de Lille (Polytech Lille)}{Villeneuve d'Ascq, France}{}{}
\cventry[.05em]{2023 à 2025}{Cycle préparatoire (MPCSI)}{\cvlogo[13pt]{polytech-mark.png}Polytech Lille}{Villeneuve d'Ascq, France}{}{Diplômé avec félicitations, major de promotion 2023/2024}
\cventry[.05em]{2023}{Baccalauréat}{France}{}{}{Mention Très Bien}

% ---------------------------------------------------------------------------
\section{Compétences techniques}
\cvitem[.05em]{Systèmes \& infra.}{Linux (NixOS, Debian, Arch, Ubuntu), KVM/QEMU, Docker, microVMs, systemd.}
\cvitem[.05em]{Config. déclarative}{Nix/NixOS, sops, builds reproductibles, gestion d'environnements.}
\cvitem[.05em]{IA locale}{llama.cpp, vLLM.}
\cvitem[.05em]{Réseau \& sécurité}{Tailscale, ZeroTier, isolation IAM, AWS (EC2, IAM, security groups), CTF (RE, web / sûreté mémoire).}
\cvitem[.05em]{Programmation}{Python, C/C++, C\#, Java, Rust, Bash, PostgreSQL, assembleur (AVR/ARM).}
\cvitem[.05em]{Outils}{Git, KiCad, MATLAB/Simulink, Matplotlib.}
\cvitem[.05em]{Langues}{anglais (C1), arabe (C2), français (C1), allemand (B2).}

% ---------------------------------------------------------------------------
\section{Projets sélectionnés}
\cventry[.05em]{Depuis août 2026}{Inférence LLM locale sur matériel grand public}{}{}{}{
\begin{cvlist}
  \item Modèles à poids ouverts sur trois machines avec \texttt{llama.cpp}, inférence GPU et CPU seul, joignables seulement par passerelle VPS : aucune machine à domicile n'ouvre de port.
  \item Un mélange d'experts de 125 Md de paramètres sur une carte de 12 Go : experts répartis entre VRAM, RAM et une table SSD de 29 Go. 50 tok/s en décodage et 1020 en préfill.
\end{cvlist}}

\cventry[.05em]{Depuis janv. 2026}{Infrastructure serveur Linux auto-hébergée}{}{}{}{
\begin{cvlist}
  \item Serveur headless : Nextcloud, Jellyfin, SearXNG, Samba. Sous Debian en janv. 2026, reconstruit sous NixOS en févr. depuis un seul arbre de configuration.
  \item Rien n'écoute sur internet : services derrière la passerelle VPS, secrets du dépôt chiffrés avec sops.
\end{cvlist}}

\cventry[.05em]{Janv. à juin 2026}{Clavier ergonomique split \textbf{(projet de cours)}}{}{}{}{
\begin{cvlist}
  \item PCB deux couches KiCad : matrice 30 touches, diodes anti-ghosting, RGB 48 canaux sur deux TLC5947. Firmware C : LUFA USB HID, balayage, lien UART vers l'autre moitié.
  \item Surintensité du rétroéclairage ayant détruit le MCU : résistance IREF du pilote corrigée et circuit remplacé ; la carte flashe et fonctionne.
\end{cvlist}}


% ---------------------------------------------------------------------------
\section{Expériences professionnelles}
\cventry[.05em]{Juin à juil. 2026}{Stage de recherche, CRIStAL \textbf{(équipe 2XS)}}{\cvlogo[13pt]{univ-lille-mark.png}IRCICA \textbullet Université de Lille}{Lille, France}{}{
\begin{cvlist}
  \item Retiré chaque fork and exec de \texttt{iw} et \texttt{ip} du module wifi d'une sonde Linux de détection d'objets connectés au profit d'appels directs \texttt{nl80211} et rtnetlink, avec découverte automatique des interfaces.
  \item La réécriture a été adoptée par la sonde, qui tourne désormais dans un seul processus au lieu d'un par appel sans-fil.
\end{cvlist}}

\cventry[.05em]{Avr. 2025 à ce jour}{Fabmanager \textbf{(FabLab Fabricarium)}}{\cvlogo[13pt]{polytech-mark.png}Polytech Lille}{Villeneuve d'Ascq, France}{}{
\begin{cvlist}
  \item Exploitation et maintenance du FabLab de l'école : imprimantes 3D, découpeuses laser, CNC et postes d'électronique.
  \item Accompagnement des prototypes matériels et IoT des étudiants et chercheurs ; réservations, sécurité et sessions ouvertes.
\end{cvlist}}

\cventry[.05em]{Juil. à août 2023}{Stage en ingénierie de l'automatisation}{\cvlogo[10pt]{valeo.png}Valeo}{Tunis, Tunisie}{}{
\begin{cvlist}
  \item Prototype de contrôle qualité par reconnaissance sonore sur une ligne de volants (Python, Librosa) ; stage d'un mois, non déployé sur la ligne.
\end{cvlist}}

\end{document}
